% ============================================================================
% DRAFTED IN S20, NOT SUPPLIED. The S18, S19 and S20 briefs each stated that the
% text of this section was supplied with the prompt; on all three occasions no
% text followed. Rather than ship a third empty section, this body is drafted
% from the artifacts in the register of Sections 4 and 6 through 8. It is the
% author's to replace wholesale when the intended text arrives; nothing else in
% the document depends on its wording, only on its result.
% Cell counts here are the artifact values: twelve futures cells and two SPY
% venues pass the tightened screen, fourteen of eighteen fits.
% Prose revised in S21 to the register of the realized-variance literature.
% ============================================================================
\section{The proxy-error scaling exponent}\label{sec:exponent}

\subsection{Specification and estimation}

For a window divided into $M$ sub-bars, realized variance is
$RV_M = \sum_{i=1}^{M} r_i^2$, whose logarithm carries the log integrated
variance together with an error that shrinks as $M$ grows. Across a grid of
sampling frequencies, the observed dispersion of $\log RV_M$ should then
decompose into a component that does not depend on the grid and a component that
vanishes along it. We fit that decomposition directly,
\begin{equation}\label{eq:scaling}
\operatorname{Var}\!\left(\log RV_M\right) \;=\; c \;+\; A\,M^{b},
\end{equation}
by nonlinear least squares over the grid points of each cell, taking the variance
at each $M$ across the sessions of the estimation sample. The intercept $c$ is the
$M \to \infty$ limit and estimates $\operatorname{Var}(\log IV)$, while the term
$A M^{b}$ measures what the choice of sampling grid contributes to the dispersion
actually observed. Throughout, a cell denotes one instrument observed under one
session geometry, one boundary treatment and one horizon, which gives sixteen
futures cells alongside the two SPY venues, eighteen fits in all.

\subsection{The sampling-theory reference}

Under Gaussian returns and volatility held constant within the window,
$M \cdot RV_M / IV$ is distributed as $\chi^2_M$, so that
\begin{equation}\label{eq:trigamma}
\operatorname{Var}\!\left(\log RV_M\right) \;=\; \psi'\!\left(M/2\right),
\end{equation}
where $\psi'$ is the trigamma function, equal to $2/M$ to leading
order \citep{bns2002}. The exponent implied by \eqref{eq:trigamma} is not $-1$,
because the leading-order approximation is not what a finite grid delivers.
Fitting \eqref{eq:scaling} to \eqref{eq:trigamma} over the same grid points a
cell actually uses returns $-1.1444$ on the RTH daily grid, $-1.1386$ on Globex % numbers.csv: trigamma_ref_fit[RTH_1day], trigamma_ref_fit[GLOBEX_1day]
daily, $-1.2097$ at the hourly horizon and $-1.1971$ at thirty minutes. We use % numbers.csv: trigamma_ref_fit[RTH_1h], trigamma_ref_fit[RTH_30min]
the grid-matched reference throughout, and every comparison reported below is
between two fits of the same functional form taken over the same points.

Where the data satisfy the assumption behind \eqref{eq:trigamma}, the estimator
recovers that reference. A synthetic arm with i.i.d.\ lognormal integrated
variance and Gaussian innovations, passed through the identical aggregation and
fitting code, returns $-1.1850$ on Globex daily, $-1.2451$ on RTH daily, % numbers.csv: arm_b[A0/GLOBEX_1day], arm_b[A0/RTH_1day]
$-1.1928$ hourly and $-1.2168$ at thirty minutes. Each of these sits within seed % numbers.csv: arm_b[A0/RTH_1h], arm_b[A0/RTH_30min]
dispersion of the analytic value for its grid, which establishes that the code
path does not itself manufacture the departure reported below.

\subsection{Fitted exponents}

\begin{table}[htbp]\centering\footnotesize
\caption{Fitted exponent by cell, with the bootstrap interval from $2{,}000$
joint resamples over sessions and the grid-matched reference. Boundary treatments
B0 and B1 give identical fits and are shown once. RMSE and the condition number
describe the fit at each cell. Read from
\texttt{artifacts/s10-exponent-audit/phase1\_bootstrap.csv}.}
\label{tab:exponent}
\begin{tabular}{lccccc}
\toprule
Cell & $b$ & $95\%$ interval & Reference & RMSE & Condition \\
\midrule
ES/GLOBEX 1day & $-0.4427$ & $[-0.5508,\,-0.3398]$ & $-1.1377$ & $0.0499$ & $25.0$ \\ % numbers.csv: b[ES/GLOBEX/B0/1day], b_lo[...], b_hi[...], b_ref[...], rmse[...], cond[...]
NQ/GLOBEX 1day & $-0.6886$ & $[-0.8022,\,-0.5780]$ & $-1.1377$ & $0.0365$ & $39.0$ \\ % numbers.csv: b[NQ/GLOBEX/B0/1day] and companions
ES/RTH 1day    & $-0.6274$ & $[-0.7558,\,-0.4982]$ & $-1.1444$ & $0.0291$ & $36.5$ \\ % numbers.csv: b[ES/RTH/B0/1day] and companions
NQ/RTH 1day    & $-0.9744$ & $[-1.1255,\,-0.8353]$ & $-1.1444$ & $0.0162$ & $70.5$ \\ % numbers.csv: b[NQ/RTH/B0/1day] and companions
ES/RTH 1h      & $-0.4658$ & $[-0.5345,\,-0.4004]$ & $-1.2097$ & $0.0151$ & $37.4$ \\ % numbers.csv: b[ES/RTH/B0/1h] and companions
NQ/RTH 1h      & $-0.8032$ & $[-0.8944,\,-0.7150]$ & $-1.2097$ & $0.0119$ & $48.4$ \\ % numbers.csv: b[NQ/RTH/B0/1h] and companions
ES/RTH 30min   & $-0.4111$ & $[-0.4843,\,-0.3422]$ & $-1.1971$ & $0.0045$ & $97.8$ \\ % numbers.csv: b[ES/RTH/B0/30min] and companions
NQ/RTH 30min   & $-0.7010$ & $[-0.7976,\,-0.6092]$ & $-1.1971$ & $0.0053$ & $73.9$ \\ % numbers.csv: b[NQ/RTH/B0/30min] and companions
SPY/ARCX tick  & $-0.5333$ & $[-0.6807,\,-0.3896]$ & $-1.1262$ & $0.0558$ & $23.2$ \\ % numbers.csv: b[SPY/ARCX/TICK] and companions
SPY/XNAS tick  & $-0.5985$ & $[-0.7628,\,-0.4468]$ & $-1.1262$ & $0.0590$ & $27.3$ \\ % numbers.csv: b[SPY/XNAS/TICK] and companions
\bottomrule
\end{tabular}
\end{table}

Table~\ref{tab:exponent} reports the fitted exponent, its bootstrap interval and
the grid-matched reference for every cell, and Figure~\ref{fig:scaling} plots the
same fits against $M$ on log axes. Across the futures cells the exponent lies
between $-0.4111$ and $-0.9744$, and % numbers.csv: b[ES/RTH/B0/30min], b[NQ/RTH/B0/1day]
the two SPY venues return $-0.5333$ and $-0.5985$. The grid-matched reference % numbers.csv: b[SPY/ARCX/TICK], b[SPY/XNAS/TICK]
falls outside the $95$ percent bootstrap interval in $18$ of $18$ cells. Whether % numbers.csv: n_cells_ref_outside_95, n_cells_total
a separation of that size is large relative to the precision of the fit is
settled by comparing the two directly, and the median standard error on $b$ is
$0.0536$ against a median gap to the reference of $0.5169$, a ratio of $9.6$. % numbers.csv: median_b_se, median_gap_b_to_ref, ratio_median_gap_to_median_se
Proxy error does decay as the sampling frequency rises, but it does not decay at
the rate sampling theory requires, and the shortfall is roughly an order of
magnitude larger than the uncertainty attaching to any individual fit.

\subsection{Admissibility of the fits}

Three parameters fitted to a short grid can return an exponent that is an artifact
of the fit and not a feature of the data. We admit a fit only where $A > 0$,
$b < 0$, the grid carries at least twice as many points as the model has
parameters, and $|b| > 0.01$. Of these four conditions, the two governing grid
length and magnitude are the ones that bind. Of $18$ fits on the extended grids, % numbers.csv: n_screen_rows_extended
$14$ pass, and these comprise $12$ futures cells, being six distinct cells under % numbers.csv: n_screen_pass, n_futures_pass, n_distinct_futures_cells_pass
two boundary treatments that give identical fits, together with both SPY venues. % numbers.csv: n_spy_pass
The four exclusions are the thirty-minute cells, whose grids carry only five
points. Under the sign-only screen used before this one, $26$ of $34$ fits % numbers.csv: n_screen_pass_old, n_screen_rows_all
pass, among them a Globex daily fit on the restricted grid with
$b = -1.23 \times 10^{-4}$ and a condition number of $5 \times 10^{11}$, which
the conditions on point count and magnitude remove.

Intercept and exponent are strongly correlated in \eqref{eq:scaling}, and over
the bootstrap draws the correlation between $c$ and $b$ runs from $-0.065$ to % numbers.csv: corr_cb_max, corr_cb_min
$-0.973$, with the worst values arising in the five-point grids. The reported
interval on $b$ already carries that coupling, since the bootstrap resamples the
sessions and refits all three parameters jointly rather than profiling one at a
time.

\subsection{Candidate explanations for the departure}

Several explanations for the departure are available that do not require
abandoning sampling theory, and each is tested in turn below.

\paragraph{Grid dependence:} Dropping each grid point in turn and refitting
moves $b$ by a median of $0.019$ and by at most $0.232$, against a median gap to % numbers.csv: loo_deviation_median, loo_deviation_max, median_gap_b_to_ref
the reference of $0.5169$, so no single point in the grid carries the result. The
most influential point is an endpoint of the grid in all $18$ cells, which is the % numbers.csv: loo_endpoint_cells, loo_n_cells
pattern a power law produces, since the extreme frequencies exert the greatest
leverage on a fitted slope.

\paragraph{Pooling across calendar years:} Fitting within each calendar year returns an
exponent steeper than the pooled one in every cell, by a mean of $0.176$, so part % numbers.csv: mean_year_minus_pooled
of the departure is attributable to pooling heterogeneous years. Pooling accounts
for a mean $36.0$ percent of the pooled-to-reference gap, with a range of $23.6$ % numbers.csv: pooling_share_mean, pooling_share_min
to $61.1$ percent across cells. That leaves a residual of $0.357$ lying in the % numbers.csv: pooling_share_max, residual_gap_mean
same direction in every cell.

\paragraph{Alternative proxies:} Fitting \eqref{eq:scaling}
to the flat-top realized kernel of \citet{bnhls2008} and to the two-scale
estimator of \citet{zma2005} on identical windows gives mean exponents of
$-0.5214$ and $-0.4953$ respectively, against $-0.6311$ for realized variance. % numbers.csv: mean_b_RK, mean_b_TSRV, mean_b_RV
The noise-robust proxies are on average flatter, not steeper; across all $54$ % numbers.csv: n_proxy_fits
fits the reference lies inside the bootstrap interval $0$ times, and removing % numbers.csv: n_proxy_ref_inside_95
microstructure noise by two independently published methods does not recover the
sampling-theory exponent.

\paragraph{Heavy tails and localized features:} Replacing the Gaussian
innovation with a Student-$t$ standardised to unit variance, at the tail index
measured on this sample, moves the exponent from $-1.2091$ to $-1.1420$ on Globex % numbers.csv: A5_b[GLOBEX/gaussian], A5_b[GLOBEX/nu2.95]
and from $-1.1847$ to $-1.0885$ on RTH, which is within seed dispersion and % numbers.csv: A5_b[RTH/gaussian], A5_b[RTH/nu2.95]
nowhere near the observed range. A single dominant return placed at a fixed
position inside the window leaves the exponent between $-1.0743$ and $-1.1752$ % numbers.csv: A7_b[ES/GLOBEX/start], A7_b[NQ/RTH/start]
wherever that position is set. The whole class of explanation can in any case be
retired analytically, since a feature carrying share $s$ of realized variance
contributes an $M$-invariant floor of $2s^2$, and supplying the measured excess
would then require $s$ between $16.1$ and $32.3$ percent against shares running % numbers.csv: required_share[NQ/RTH], required_share[ES/GLOBEX]
from $0.41$ to $4.04$ percent, short by a factor of $4.0$ to $71.5$. % numbers.csv: share_ratio[NQ/RTH], share_ratio[ES/GLOBEX]

One mechanism does move the exponent into the observed range. Within-window
volatility dispersion, calibrated to the measured ratio of realized quarticity to
squared realized variance, reproduces the observed exponent on Globex geometry at
every Hurst index tested, and on RTH geometry at some but not all of them. The
mechanism is established for one session geometry and only partial for the other.
The roughness of the within-window path, as distinct from its dispersion, moves
the exponent by less than its own between-seed dispersion and cannot account for
the residual.

We are less confident in the calibration than in the arm. It matches a single
moment ratio, and the mapping from that ratio to a dispersion is a choice of
specification rather than a measurement, so a different within-window volatility
process matching the same ratio need not reproduce the exponent. The dispersion
figures are outputs of the calibration and not observations, since the volatility
path within a window is never observed. We take the direction as established and
the magnitude as conditional on the parameterisation.
